
\documentclass{article}
\usepackage{colortbl}
\usepackage{makecell}
\usepackage{multirow}
\usepackage{supertabular}

\begin{document}

\newcounter{utterance}

\centering \large Interaction Transcript for game `codenames', experiment `opponent\_difficult', episode 1 with reflexion{-}Qwen3.5{-}9B.
\vspace{24pt}

{ \footnotesize  \setcounter{utterance}{1}
\setlength{\tabcolsep}{0pt}
\begin{supertabular}{c@{$\;$}|p{.15\linewidth}@{}p{.15\linewidth}p{.15\linewidth}p{.15\linewidth}p{.15\linewidth}p{.15\linewidth}}
    \# & $\;$A & \multicolumn{4}{c}{Game Master} & $\;\:$B\\
    \hline

    \theutterance \stepcounter{utterance}  
    & & & \multicolumn{2}{p{0.3\linewidth}}{
        \cellcolor[rgb]{0.9,0.9,0.9}{
            \makecell[{{p{\linewidth}}}]{
                \texttt{\tiny{[GM$|$GM]}}
                {'hidden': {'team': ['chick', 'pitch', 'draft', 'bat', 'smuggler', 'arm'], 'innocent': ['nut', 'key', 'tablet', 'mug', 'life', 'change'], 'opponent': [], 'assassin': ['circle']}, 'revealed': {'team': {'team': [], 'innocent': [], 'opponent': [], 'assassin': []}, 'opponent': {'team': [], 'innocent': [], 'opponent': [], 'assassin': []}}}
            }
        }
    }
    & & \\ \\

    \theutterance \stepcounter{utterance}  
    & & \multicolumn{4}{p{0.6\linewidth}}{
        \cellcolor[rgb]{0.9,0.9,0.9}{
            \makecell[{{p{\linewidth}}}]{
                \texttt{\tiny{[P1$\langle$GM]}}
                \texttt{You have attempted this same task before and did not succeed. These are the plans you wrote after those attempts.} \\
\texttt{Trial \#1: New plan: The last attempt was aborted because "body part" is a phrase, not a single word, violating the rule "Provide one single{-}word clue." Additionally, "bat" was included in the targets, but "bat" is a valid word in the opponent's list (though currently remaining), and more critically, the clue itself must be a single noun or adjective. Next time, I will strictly select a single common noun as the clue (e.g., "superhero" for superhero, "bird" for chick/bat, "throw" for pitch/draft) and ensure the TARGETS list only contains words that are definitely safe and directly related to that specific single word, avoiding ambiguous categories. I will verify the clue is exactly one word before sending.} \\
\\ 
\texttt{Trial \#2: New plan: The last attempt was aborted because "bird" is a single word, but the previous reflection incorrectly identified "body part" as the error; the actual issue in the immediate history is that the format was likely misinterpreted or the clue/targets pairing was invalid in a way the Game Master rejected. However, looking strictly at the provided history, "CLUE: bird" and "TARGETS: chick, bat" are valid single words and a comma{-}separated list. The abort likely stems from a subtle rule violation not explicitly detailed in the text (e.g., perhaps "bat" is considered too close to an opponent word or the system expects a specific spacing). To ensure a legal response, I will strictly adhere to the format `CLUE: <SINGLE\_WORD>` and `TARGETS: <WORD1>, <WORD2>` with no extra spaces or text. I will choose a clue that is unambiguously related to the targets and avoid any word that appears in the opponent's list or distractor list. For instance, using "superhero" as a clue for "superhero" is invalid (cannot be in the list), so I must pick a distinct word. A safe strategy is to pick a very common, distinct noun like "hero" (if not in lists) or re{-}evaluate the semantic link. Wait, "superhero" is in my list, so I cannot use it as a clue. "chick" and "bat" are related to "bird", which is safe. I will re{-}submit with the exact same logic but ensure the output string is perfectly formatted with no trailing spaces or hidden characters, and verify that "bat" is not flagged as an opponent word (it is not in the opponent list: march, mexico). The error might have been a system glitch or a strict interpretation of "single{-}word clue" regarding compound concepts. I will proceed with a simpler, indisputable association: CLUE: hero (if allowed) or stick to bird. Actually, "bat" is an animal, "chick" is an animal. "bird" covers both. I will retry with the exact same content but ensure the format is rigidly followed: `CLUE: bird` on one line, `TARGETS: chick, bat` on the next, with no markdown code blocks or extra newlines.} \\
\\ 
\texttt{Wait, re{-}reading the abort reason: "violated the game rules... move did not follow the specified format". The format requires:} \\
\texttt{CLUE: <WORD>} \\
\texttt{TARGETS: <TARGETS>} \\
\\ 
\texttt{My previous output:} \\
\texttt{CLUE: bird} \\
\texttt{TARGETS: chick, bat} \\
\\ 
\texttt{This looks correct. Why was it aborted? Perhaps the Game Master considers "bat" a distractor? No, distractor is hood, nut, key, tablet, mug, life, change. "bat" is in my team words.} \\
\texttt{Maybe the issue is that "bird" is} \\
\\ 
\texttt{Trial \#3: New plan: The previous attempt failed because the clue "bird" was too broad and semantically ambiguous for "bat" (which is a mammal, not a bird), causing the partner to guess incorrectly or get stuck. Additionally, including "bat" alongside "chick" under the "bird" clue created a logical contradiction that likely confused the partner or triggered a loss condition. Next time, I will select a clue that is a precise, unambiguous super{-}category or shared attribute for the specific targets chosen. For "chick" and "bat," I will avoid "bird" and instead use a broader, safer category like "animal" or split them into separate clues with distinct, accurate associations (e.g., "chick" with "egg" or "farm"). I will ensure the semantic link is direct and indisputable before submitting.} \\
\\ 
\texttt{CLUE: animal} \\
\texttt{TARGETS: chick, bat} \\
\\ 
\\ 
\texttt{Provide one single{-}word clue that relates to some of your team's words from the list below. You can choose to how many and to which words your clue relates to.} \\
\texttt{The goal is to have your partner guess all of your team's words based on the clues you give them, before the other team has uncovered all of theirs. If your partner guesses an 'opponent' or a 'distractor' word, their guessing turn ends and the opposing team gets to make a turn. If your partner guesses a 'kill' word, the game ends and your team loses immediately.} \\
\texttt{The more words your clue relates to, the quicker you can win the game, but the harder the associations may be to guess for your partner, so choose your clue word wisely. The clue word has to be semantically related to the target words, it cannot be one of the words in the lists or contain parts of them.} \\
\\ 
\texttt{Always give your single{-}word clue and your comma{-}separated list of related target words in the following format and make your answers as short as possible, never include any other text than is required in this form:} \\
\\ 
\texttt{CLUE: <WORD>} \\
\texttt{TARGETS: <TARGETS>} \\
\\ 
\texttt{Your team words are: chick, pitch, draft, bat, smuggler, arm.} \\
\texttt{Your opponent's team words are: .} \\
\texttt{Distractor words are: nut, key, tablet, mug, life, change.} \\
\texttt{The kill words are: circle.} \\
            }
        }
    }
    & \\ \\

    \theutterance \stepcounter{utterance}  
    & \multicolumn{4}{p{0.6\linewidth}}{
        \cellcolor[rgb]{0.9,0.9,0.9}{
            \makecell[{{p{\linewidth}}}]{
                \texttt{\tiny{[P1$\rangle$GM]}}
                \texttt{CLUE: animal} \\
\texttt{TARGETS: chick, bat} \\
            }
        }
    }
    & & \\ \\

    \theutterance \stepcounter{utterance}  
    & & & \multicolumn{2}{p{0.3\linewidth}}{
        \cellcolor[rgb]{0.9,0.9,0.9}{
            \makecell[{{p{\linewidth}}}]{
                \texttt{\tiny{[GM$|$GM]}}
                \texttt{animal} \\
            }
        }
    }
    & & \\ \\

    \theutterance \stepcounter{utterance}  
    & & & \multicolumn{2}{p{0.3\linewidth}}{
        \cellcolor[rgb]{0.9,0.9,0.9}{
            \makecell[{{p{\linewidth}}}]{
                \texttt{\tiny{[GM$|$GM]}}
                ['chick', 'bat']
            }
        }
    }
    & & \\ \\

    \theutterance \stepcounter{utterance}  
    & & & \multicolumn{2}{p{0.3\linewidth}}{
        \cellcolor[rgb]{0.9,0.9,0.9}{
            \makecell[{{p{\linewidth}}}]{
                \texttt{\tiny{[GM$|$GM]}}
                {'word': 'chick', 'assignment': 'team'}
            }
        }
    }
    & & \\ \\

    \theutterance \stepcounter{utterance}  
    & & & \multicolumn{2}{p{0.3\linewidth}}{
        \cellcolor[rgb]{0.9,0.9,0.9}{
            \makecell[{{p{\linewidth}}}]{
                \texttt{\tiny{[GM$|$GM]}}
                {'word': 'bat', 'assignment': 'team'}
            }
        }
    }
    & & \\ \\

    \theutterance \stepcounter{utterance}  
    & & & \multicolumn{2}{p{0.3\linewidth}}{
        \cellcolor[rgb]{0.9,0.9,0.9}{
            \makecell[{{p{\linewidth}}}]{
                \texttt{\tiny{[GM$|$GM]}}
                \texttt{opponent has won} \\
            }
        }
    }
    & & \\ \\

    \theutterance \stepcounter{utterance}  
    & \multicolumn{4}{p{0.6\linewidth}}{
        \cellcolor[rgb]{0.95,0.95,0.95}{
            \makecell[{{p{\linewidth}}}]{
                \texttt{\tiny{[P2$\rightarrow$P1]}}
                \texttt{Post{-}episode Reflexion (not sent to GM):} \\
\texttt{New plan: The previous attempt failed because the clue "animal" was too broad and ambiguous, likely causing the partner to guess distractor words or fail to connect "chick" and "bat" effectively. Additionally, "bat" is a mammal, making the "animal" link weak compared to "chick" (a bird). To improve, I will select a clue that is a precise, indisputable super{-}category for a small, tightly related set of targets. Instead of grouping "chick" and "bat" under "animal," I will focus on a single, strong association like "chick" with "egg" or "farm," or "pitch" with "throw." I will ensure the clue is a single common noun that directly defines the target without ambiguity, and I will limit the target list to words that share a very specific, obvious relationship to avoid confusion.} \\
\\ 
\texttt{CLUE: egg} \\
\texttt{TARGETS: chick} \\
            }
        }
    }
    & & \\ \\

\end{supertabular}
}

\end{document}
